\section{Related Work}
\label{sec:related}

\paragraph{Variance Decomposition in ML/NLP Evaluation.}
Factorial decomposition of evaluation variance has roots in the ML reproducibility literature. \citet{bouthillier2021accounting} introduced an ANOVA-based template that decomposes benchmark variance across implementation details, hyperparameters, and data splits, establishing that design choices account for a substantial share of reproducibility failures. Several subsequent studies applied similar decomposition logic to specific facets of model evaluation \citep{dodge2020finetuning, henderson2018deep, magnusson2023reproducibility}. \citet{mesting2026hidden} brought Generalizability theory into LLM evaluation, decomposing variance across five methodology facets on a single benchmark; judge identity dominated at 43.8\%, with prompt$\times$judge interaction at 12.4\%. However, this study used a single benchmark, included no infrastructure facets such as numerical precision, and fixed the judge as an LLM evaluator rather than isolating other sources. \citet{kunievsky2025intent} proposed an alternative decomposition that separates intent uncertainty, articulation variance, and model uncertainty, but this framework does not cross infrastructure and methodology facets in a factorial design. Our work extends these approaches by crossing seven facets across four benchmarks, enabling direct comparison of variance structures across evaluation formats.

\paragraph{Infrastructure and Prompt Variance.}
A parallel line of work has examined individual variance sources in isolation. On the infrastructure side, \citet{yuan2025layercast} documented up to 9\% accuracy variation from numerical precision and GPU type but provided no factorial attribution. \citet{messina2026background} formalized the concept of background temperature from floating-point non-associativity in GPU computation, and \citet{thinkingmachines2025batch} proposed batch-invariant kernels to mitigate CUDA nondeterminism. On the methodology side, \citet{sclar2024formatting} found up to 76\% accuracy swings from prompt formatting changes on fixed hardware, and BrittleBench \citep{brittlebench2026} reported 50\% variance from prompt perturbation. \citet{pezeshkpour2024ordering} showed that answer ordering in multiple-choice benchmarks introduces systematic bias, while \citet{bui2025seeds} and \citet{biderman2024lessons} documented sensitivity to random seeds and training configurations. Each of these studies isolates one factor while holding others fixed, which precludes estimating interactions between sources or comparing their relative magnitudes within a unified design.

\paragraph{IRT and Efficient Benchmarking.}
Item Response Theory (IRT) offers a complementary perspective on evaluation design by modeling item difficulty and discrimination parameters. \citet{lalor2019irt} introduced IRT-based item analysis for NLP benchmarks, and \citet{rodriguez2021leaderboards} showed that IRT reveals hidden properties of evaluation items that raw accuracy scores obscure. Building on this foundation, \citet{polo2024tinybenchmarks} developed tinyBenchmarks, using IRT-based subsampling to reduce evaluation cost while preserving ranking accuracy, and \citet{zhou2026irt} recommended a 116-item subset via PSN-IRT for efficient MMLU evaluation. \citet{perlitz2023efficient} explored related strategies for benchmark compression. These approaches optimize item selection for latent trait estimation but do not model variance from prompt, seed, or ordering facets. IRT-selected items evaluated under a single methodological condition can therefore systematically underestimate the number of items required for reliable measurement across diverse evaluation settings.

\paragraph{Binary Data and Psychometrics Foundations.}
Our application of G-theory to binary correctness scores draws on established psychometric methodology. \citet{shavelson1991generalizability} and \citet{brennan2001generalizability} provide the canonical framework for multi-facet reliability analysis. \citet{mushquash2006gtheory} extended G-theory to binary data, and \citet{jiang2018bayesian} developed logistic-link generalized G-theory for binary outcomes. \citet{tuerlinckx2006glmm} established the connection between generalized linear mixed models and G-theory variance components. These methods ground our use of Henderson Method~I for binary correctness scores (\S\ref{sec:binary}).

\paragraph{Positioning.}
Relative to \citet{mesting2026hidden}, we fix the judge (generation-based scoring) to isolate the remaining facets; their judge-dominated variance (43.8\%) would otherwise mask all other sources. Our multi-benchmark design reveals format-conditional structure that single-benchmark studies cannot detect. Relative to IRT approaches \citep{polo2024tinybenchmarks, zhou2026irt}, G-theory and IRT optimize orthogonal dimensions of evaluation design: IRT selects informative items, while G-theory determines the conditions needed for a target reliability. Our D-study quantifies this complementarity: multi-facet designs require $8.8\times$ fewer items than single-condition protocols for equivalent reliability. Relative to isolated-factor studies \citep{yuan2025layercast, sclar2024formatting, brittlebench2026}, our design crosses infrastructure and methodology facets within a single factorial framework, enabling interaction estimation that per-factor studies cannot provide.
